% SPDX-License-Identifier: CC-BY-SA-4.0
% Session: Form parallelism to concurrency
% Exercise: Parallelization of matrix multiplication

\begingroup

\begin{exercice}[Parallélisation]
  \label{exo:introduction/matrix}

  Le but de cet exercice est d'écrire une fonction \lstinline{double[][] product(double[][] M1, double[][] M2)}
  qui calcule le produit de deux matrices carrées $M_1$ et $M_2$ de taille $n\times n$.
  On rappelle la définition de la multiplication des matrices. Si $M_3 = M_1 \times M_2$, alors pour tout
  $i,j \in \{0,\dots,n-1\}$,
  $$ M_3[i][j] = \sum_{k=0}^{n-1} M_1[i][k] \times M_2[k][j].$$

  \begin{question}
  \item Implémentez la version séquentielle de \lstinline{product}.
  \end{question}
  \begin{correction}
    \begin{lstlisting}[gobble=6]
      public interface Product {
	double[][] product(double[][] M1, double[][] M2) throws InterruptedException;
      }
      public final class Sequential implements Product {
	double[][] product(double[][] M1, double[][] M2) {
	  int n = M1.length;
	  double[][] M3 = new double[n][n];
	  for (int i = 0; i < n; i++) {
	    for (int j = 0; j < n; j++) {
	      M3[ii][jj] = 0;
	      for (int k = 0; k < n; k++) {
		M3[ii][jj] += M1[ii][k] * M2[k][jj];
	      }
	    }
	  }
          return M3;
        }
      }
    \end{lstlisting}
  \end{correction}

  \begin{question}
  \item Proposez une parallélisation où chaque cellule $M_3[i][j]$ est calculée par un thread distinct (jusqu'à $n^2$ threads). 
  \end{question}
  \begin{correction}
    \begin{lstlisting}[gobble=6]
      package td1_parallelism.matrix;

      public final class Naive implements Product {

	public double[][] product(double[][] M1, double[][] M2) throws InterruptedException {

	  int n = M1.length;
	  double[][] M3 = new double[n][n];
	  Thread[] threads = new Thread[n*n];

	  for (int i = 0; i < n; i++) {
	    for (int j = 0; j < n; j++) {
	      final int ii = i, jj = j;
	      threads[i*n+j] = new Thread(() -> {
		M3[ii][jj] = 0;
		for (int k = 0; k < n; k++) {
		  M3[ii][jj] += M1[ii][k] * M2[k][jj];
		}
	      });
	    }
	  }

	  for(var t : threads) t.start();
	  for(var t : threads) t.join();

	  return M3;
	}

      }
    \end{lstlisting}
  \end{correction}

  \begin{question}
  \item Proposez une version qui crée au plus $T$ threads ($T$ est une variable globale), chacun calculant un lot de cellules.
    Précisez votre stratégie de répartition.
  \end{question}
  \begin{correction}
    \begin{lstlisting}[gobble=6]
      public static double[][] product(int n, double[][] M1, double[][] M2) {
        double[][] M = new double[n][n];
        Thread threads[][] = new Thread[n][n];
        for(int i = 0; i < n; i++) {
	  int thisI = i;
	  threads[i][0] = new Thread(() -> {
	    for(int j = 0; j < n; j++) {
	      int thisJ = j;
	      threads[thisI][j] = new Thread(() -> {
	        M[thisI][thisJ] = 0;
	        for(int k = 0; k<n; k++) {
		  M[thisI][thisJ]+=M1[thisI][k] * M2[k][thisJ];
	        }
	      });
	      threads[thisI][j].start();
	    }
	    for(int j = 1; j < n; j++) {
	      try {
                threads[thisI][j].join();
              } catch (InterruptedException e) { }
	    }
	    
	  });
	  threads[i][0].start();
        }
        for(int i = 1; i < n; i++) {
	  try {
            threads[i][0].join();
          } catch (InterruptedException e) { }
        }
        return M;
      }
    \end{lstlisting}
  \end{correction}

  \begin{question}
  \item Analysez la complexité en temps des versions et justifiez l'intérêt
    de limiter le nombre de threads à une borne indépendante de $n$.
  \end{question}
  \begin{correction}
    Si on parallélise la boucle intérieure, on passe d'une complexité $O(n^3)$ à $O(n^2)$, soit une accélération en $O(n)$.
    On peut faire mieux en parallélisant la boucle extérieure également. On passe alors de $O(n^3)$ à $O(n)$, soit une accélération en $O(n^2)$.
    Cela ne viole pas la loi d'Amdahl car on change le problème en changeant le nombre de threads :
    on augmente le nombre de threads pour s'adapter à la quantité de données à traîter. C'est l'idée de la loi de Gustafson :
    Si le nombre de threads est proportionnel à la taille des données, on peut espérer une accélération en $O(n)$.
    
    Remarque : le programme ci-dessus utilise $n^2 + n + 1$ threads, mais on peut utiliser les mêmes threads pour paralléliser et pour faire certains calculs.
  \end{correction}
  
\end{exercice}

\endgroup
\endinput
